\chapter{Introduction\label{chap:introduction}}
% 中文版思路：
% 在人工智能和机器人领域，工程实施是数学模型、算法创新和实际生产力提升之间的关键桥梁。当复杂的数学公式和人工智能理论框架通过工程方法有效地转化为实际产品时，其推动人类生产力发展的潜力将得到更充分的体现。多年来，机器人领域出现了各种版本的定位与地图构建技术，比如：Gmapping、Cartographer、ORB-SLAM、FAST-LIO、VINS-Mono、VINS-Fusion等。 
% 本论文阐述了一个将基于双目相机的纯视觉SLAM技术工程化为室内递送系统产品的过程。本阶段，该系统产品的功能包括3个方面：
% 1. 室内机器人进行同时定位与建图；
% 2. 建图完成后，机器人在室内行走时，可以在地图上实时反映出机器人的位置；
% 3. 设定一个目标位置，当机器人到达目标位置时，无人机可以自动起飞运送货物到目的地；
% 未来工作：
% 1.设定一个目标位置，机器人可以自动导航至目的地；
% 2.将该送货系统的工作场景，扩展至室外。
% 该工程使用的硬件设备主要有：Turtlebot3 Waffle Pi, ZED 2 双目相机，边缘计算单元Jetson Orin Nano (Max power 15W), Battery: INIU 65W 20000mAh Power Bank, 大疆Tello EDU版无人机；监控系统使用一台一般的内存16GB的笔记本电脑或PC机即可。软件模块主要有Ubuntu 22.04 LTS Server (用于Raspberry Pi SBC的操作系统), 本文接下来将从系统选型、系统设计、系统实现、工程评估、结论等方面对该工程进行详细阐述。

% English version:
Engineering can enable acdamic papers to truely change the world. When sophisticated mathematical formulations and AI theoretical frameworks are effectively translated into tangible products through engineering approaches, their potential to drive human productivity reaches its fullest expression.

This thesis presents an engineering perspective on the implementation of a robotic pure vision SLAM (Simultaneous Localization and Mapping) system designed for lightweight edge computing embedded systems, which involves:

\begin{enumerate}
    \item Theoretical foundations (rtabmap and ROS2)
    \item Engineering Design
    \item Engineering Implementation
    \item Engineering Evaluation
    \item Conclusions
\end{enumerate}

The hardware utilized includes Turtlebot3 Waffle Pi, ZED 2 stereo camera, Jetson Orin Nano edge computing unit, INIU 65W 20000mAh power bank, and DJI Tello EDU drone.

The main software packages used are Ubuntu 22.04 LTS Server, Ubuntu 22.04 LTS Desktop, JetPack 6.2.1, ROS2 Humble, turtlebot3\_bringup, DynamixelSDK, rtabmap, rtabmap\_ros, zed\-ros2\-wrapper, turtlebot3\_teleop, DJITelloPy, navigation2, rviz2, rtabmap\_viz, rqt.

The conclusion is that pure visual SLAM based on binocular cameras can be applied to edge computing embedded systems with limited computing resources through engineering methods to achieve mapping, positioning, multi-agent (robot and drone) collaboration, and navigation. 